\documentclass[11pt]{article}

\usepackage[T1]{fontenc}
\usepackage{lmodern}
\usepackage[margin=1in]{geometry}
\usepackage{amsmath,amssymb,mathtools,amsthm}
\usepackage{booktabs}
\usepackage{enumitem}
\usepackage{xcolor}
\usepackage{microtype}
\usepackage[hidelinks]{hyperref}
\usepackage{url}

\allowdisplaybreaks[2]
\emergencystretch=2em

\newtheorem{theorem}{Theorem}[section]
\newtheorem{proposition}[theorem]{Proposition}
\newtheorem{lemma}[theorem]{Lemma}
\newtheorem{corollary}[theorem]{Corollary}
\newtheorem{remark}[theorem]{Remark}

\DeclareMathOperator{\tr}{tr}
\DeclareMathOperator{\rank}{rank}
\DeclareMathOperator{\Span}{span}
\newcommand{\Frob}{\mathrm F}
\newcommand{\R}{\mathbb R}
\newcommand{\cD}{\mathcal D}
\newcommand{\cF}{\mathcal F}
\newcommand{\MT}{\mathrm{MT}}
\newcommand{\Nsimple}{N_0^s}

\hypersetup{
  pdftitle={A Computer-Assisted Seven-Point Refinement for Simple Zeros of the Riemann Zeta Function},
  pdfauthor={Anonymous submission; AI-assisted research draft}
}

\title{A Computer-Assisted Seven-Point Refinement\\
for Simple Zeros of the Riemann Zeta Function}
\author{Anonymous submission\\\small AI-assisted research draft}
\date{Peer-review draft, August 14, 2026}

\begin{document}

\maketitle

\begin{center}
\fcolorbox{black}{gray!8}{\parbox{0.91\linewidth}{
\textbf{Status.} This manuscript presents an unreviewed,
computer-assisted proof of a refinement of Claude's Theorem~D.  It does not
prove the Riemann hypothesis.  The finite seven-point inequality has been
replayed locally at two Arb precision settings, and its delicate
strong-convexity basin has also been checked by a separate exact-rational LDL
implementation using the same Arb/FLINT transcendental trust base.  The
endpoint deletion, finite-grid tail, uniform Gram-to-kernel comparison, and
summation of all block errors are proved explicitly below.  The prime-side
trace asymptotics are cited from Claude's Theorem~D~\cite{Claude26}.  A
partial Lean~4 companion checks the finite-dimensional algebra but is not
used as a substitute for the analytic proof.  The result has not yet
received independent specialist peer review.}}
\end{center}

\begin{abstract}
Let $N(T,2T)$ count the nontrivial zeros of the Riemann zeta function with
ordinates in $(T,2T]$, with multiplicity, and let $\Nsimple(T,2T)$ count the
simple zeros on the critical line in the same interval.  We retain a convex
spectral defect in the rank--trace argument underlying Claude's 2026
Montgomery--Taylor bound and connect that defect to short-range correlations
of consecutive critical-line zeros.  The only new computer-assisted input is
the global inequality
\[
 \cF_6(g_1,\ldots,g_6)
 \ge C_*:=\frac{38262312113}{10^{13}}
 \qquad(g_1,\ldots,g_6\ge0),
\]
for an explicit seven-point functional built from the normalized
Montgomery--Taylor overlap kernel.  An exhaustive interval verifier uses Arb
enclosures, certified convex tangent bounds, and two rational
strong-convexity basins; a higher-precision replay has identical logical
counts.  A Lean companion formalizes the new matrix and exact-arithmetic
layers and exact final arithmetic.  We prove directly that endpoint deletion
costs $o(N)$, that the finite simple-zero Gram entries differ uniformly from
the limiting kernel by $O(1/\log T)$ on every relevant block, and that the
error summed over all 267 shifts is $o(N)$.  The Lean companion is supporting
evidence rather than a premise of the proof.  The resulting bound is
\[
 \liminf_{T\to\infty}\frac{\Nsimple(T,2T)}{N(T,2T)}
 \ge 0.6730254768378743181159739993\ldots .
\]
\end{abstract}

\section{Scope and statement of the result}

Write
\[
 N:=N(T,2T),\qquad S:=\Nsimple(T,2T),
\]
and let
\begin{equation}\label{eq:HMT}
 H_{\MT}:=\frac32-\frac1{\sqrt2}\cot\frac1{\sqrt2}
 =0.6725007036794116457343797908\ldots .
\end{equation}
The quantity $H_{\MT}$ is the optimized constant in Theorem~D of
Claude~\cite{Claude26}.  That theorem is accompanied by a Lean~4 artifact for
its stated results~\cite{ClaudeLean26}.  We use its analytic trace and tail
estimates in the ordinary manner of a cited theorem.  The accompanying
\texttt{ZetaSeven} package is a pinned partial formalization of the new
finite-dimensional layers; the mathematical proof in this manuscript does
not depend on an end-to-end Lean theorem.

The exact local constant certified in this work is
\begin{equation}\label{eq:Cstar}
 C_*:=\frac{38262312113}{10000000000000}
 =0.0038262312113.
\end{equation}

\begin{theorem}[Computer-assisted refinement]\label{thm:main}
One has
\begin{align}
 \liminf_{T\to\infty}\frac{\Nsimple(T,2T)}{N(T,2T)}
 &\ge \frac{267H_{\MT}-266/500}{267-261C_*}\label{eq:main-formula}\\
 &=\frac{2670000000000000H_{\MT}-5320000000000}
 {2660013536538507}\notag\\
 &=0.67302547683787431811597399932145\ldots .\notag
\end{align}
Equivalently, the lower proportion is
$67.3025476837874318115\ldots\%$.  Relative to the cited
Montgomery--Taylor value $H_{\MT}=67.2500703679411645\ldots\%$, the increase
is $0.0524773158462672\ldots$ percentage points.
\end{theorem}

The theorem is unconditional: no Riemann hypothesis or zero-density
hypothesis is assumed.  As with any result built on a recent preprint, its
status remains subject to checking of the cited source and independent
review of the new argument.  The adjective ``computer-assisted'' refers only
to Proposition~\ref{prop:F6} and its explicit software trust base.

The proof has three new layers.  Section~\ref{sec:stability} preserves a
nonnegative spectral defect in the matrix inequality.  Sections
\ref{sec:kernel}--\ref{sec:certificate} lower-bound that defect using the
Montgomery--Taylor kernel and a rigorous seven-point computation.  Sections
\ref{sec:block}--\ref{sec:pinching} convert the local inequality to the exact
constant in Theorem~\ref{thm:main}.

\section{The imported analytic interface}\label{sec:interface}

We state precisely the portion of~\cite{Claude26} used by the refinement.
Put
\[
 \ell:=\log\frac{T}{2\pi},\qquad L:=\ell,\qquad
 I:=(T,2T],\qquad I':=(T-T^{1/2},2T+T^{1/2}].
\]
Fix the ramp profile $\varrho$ of~\cite[Section~2.2]{Claude26} and use ramp
width one.  The optimized window of~\cite[Section~7.1]{Claude26} is
\begin{equation}\label{eq:optimized-window}
 \phi_L(u):=\cos\!\left(\frac{\sqrt2u}{L}\right)^{1/2}
 \varrho(L/2-|u|),
 \qquad
 a_L:=\frac1L\int_{\R}\phi_L(u)^2\,du.
\end{equation}
Put $h=2\pi/L$, $d=\lfloor T/h\rfloor$, and $\tau_k=T+kh$.  Let
$\widehat A=A/(a_LL^2)$ be the normalized finite zero-side compression of
Weil's Hermitian form from~\cite[Section~4]{Claude26}.
Partition the zeros in $I'$ into simple critical-line zeros $\mathcal S_1$,
distinct multiple critical-line zeros $\mathcal S_2$, and unordered off-line
pairs $\mathcal P$.  Write
\[
 s_1:=\#\mathcal S_1,\qquad s_2:=\#\mathcal S_2,
 \qquad p:=\#\mathcal P.
\]
Counting multiplicity gives
\begin{equation}\label{eq:Nlower}
 N(I')\ge s_1+2s_2+2p.
\end{equation}

For $\rho\in\mathcal S_1$, with ordinate $\gamma_\rho$, put
\[
 u_\rho:=\bigl(\widehat\phi_L(\gamma_\rho-\tau_k)\bigr)_{0\le k<d},
 \qquad v_\rho:=\frac{u_\rho}{\sqrt{a_L}\,L},
\]
and let $V$ have these atoms as columns.  The full-grid identity
\cite[Lemma~2.2]{Claude26} gives
\[
 \sum_{k\in\mathbb Z}
 |\widehat\phi_L(\gamma_\rho-\tau_k)|^2=a_LL^2,
\]
so every finite column $v_\rho$ has norm at most one.  Moreover
\[
 P_1=VV^*,\qquad M=V^*V,
\]
and for $Q':=\widehat A-P_1$ the inertia decomposition in the proof of
\cite[Proposition~4.4]{Claude26} gives
\begin{equation}\label{eq:positive-index}
 n_+(Q')\le s_2+p.
\end{equation}
The trace and tail analysis in~\cite[Sections~4--7]{Claude26} yields the
baseline constant~\eqref{eq:HMT}.  Corollary~\ref{cor:global-defect} below
shows that the same estimates, after applying the strengthened matrix
inequality, give
\begin{equation}\label{eq:analytic-interface}
 S\ge H_{\MT}N+\cD(M^\circ)-o(N),
\end{equation}
where $M^\circ$ is the Gram matrix of the retained central simple zeros and
$\cD$ is defined in~\eqref{eq:Psi} below.

The full-grid identity just quoted is exact.  Its finite-grid truncation and
the limiting kernel comparison are proved with explicit uniform error in
Lemma~\ref{lem:kernel-limit}; they are not assumed as an additional
interface.

\section{A stability refinement of the rank--trace inequality}
\label{sec:stability}

Define the convex, nonnegative function
\begin{equation}\label{eq:Psi}
 \Psi(t):=
 \begin{cases}
  (t-1)^2,&0\le t\le2,\\
  2t-3,&t\ge2,
 \end{cases}
 \qquad
 \cD(M):=\tr\Psi(M)
\end{equation}
for positive semidefinite $M$.

\begin{lemma}[Stability-enhanced rank--trace inequality]
\label{lem:stable-rank-trace}
Let $V$ be a $d\times r$ matrix whose columns have norm at most one.  Put
$P=VV^*\succeq0$ and $M=V^*V$, and let $Q$ be Hermitian with
$n_+(Q)\le b$.  Then
\begin{equation}\label{eq:stable-rank-trace}
 \|P+Q\|_{\Frob}^2
 \ge4\tr(P+Q)-3r-4b+\cD(M).
\end{equation}
\end{lemma}

\begin{proof}
Write $Q=Q_+-Q_-$ with $Q_\pm\succeq0$, $Q_+Q_-=0$, and
$\rank Q_+\le b$.  Since $P,Q_+\succeq0$,
\[
 \|P+Q\|_{\Frob}^2
 \ge\|P-Q_-\|_{\Frob}^2+\|Q_+\|_{\Frob}^2.
\]
For every $q\ge0$, $q^2\ge4q-4$, hence
\begin{equation}\label{eq:qplus}
 \|Q_+\|_{\Frob}^2\ge4\tr Q_+-4b.
\end{equation}

Let $p_1\ge\cdots\ge p_r\ge0$ be the eigenvalues of $M$, padding the
nonzero spectrum of $P$ by zeros, and let $n_1\ge n_2\ge\cdots\ge0$ be the
eigenvalues of $Q_-$.  Von Neumann's trace inequality gives
$\tr(PQ_-)\le\sum_{i\le r}p_in_i$.  Therefore
\[
 \|P-Q_-\|_{\Frob}^2+4\tr Q_-
 \ge\sum_{i\le r}\bigl((p_i-n_i)^2+4n_i\bigr).
\]
For $p\ge0$,
\begin{equation}\label{eq:scalar-min}
 \min_{n\ge0}\bigl((p-n)^2+4n\bigr)
 =\begin{cases}p^2,&0\le p\le2,\\4p-4,&p\ge2,
 \end{cases}
 =2p-1+\Psi(p).
\end{equation}
It follows that
\begin{equation}\label{eq:qminus}
 \|P-Q_-\|_{\Frob}^2
 \ge2\tr P-r+\cD(M)-4\tr Q_-.
\end{equation}
Adding~\eqref{eq:qplus} and~\eqref{eq:qminus} gives
\[
 \|P+Q\|_{\Frob}^2
 \ge4\tr(P+Q)-2\tr P-r-4b+\cD(M).
\]
Finally $\tr P$ is the sum of the squared column norms of $V$, so
$\tr P\le r$.
\end{proof}

\begin{corollary}[Global defect inequality]\label{cor:global-defect}
With the notation of Section~\ref{sec:interface},
\begin{equation}\label{eq:s1-finite}
 s_1\ge4\tr\widehat A-\|\widehat A\|_{\Frob}^2
 -2N(I')+\cD(M).
\end{equation}
After the endpoint truncation and analytic estimates of~\cite{Claude26},
\begin{equation}\label{eq:global-defect}
 S\ge H_{\MT}N+\cD(M^\circ)-o(N).
\end{equation}
\end{corollary}

\begin{proof}
Apply Lemma~\ref{lem:stable-rank-trace} to
$P_1+Q'=\widehat A$.  Equations~\eqref{eq:Nlower} and
\eqref{eq:positive-index} imply
\begin{align*}
 \|\widehat A\|_{\Frob}^2
 &\ge4\tr\widehat A-3s_1-4s_2-4p+\cD(M)\\
 &\ge4\tr\widehat A-s_1-2N(I')+\cD(M),
\end{align*}
which is~\eqref{eq:s1-finite}.

Let $M^\circ$ be the principal block belonging to the simple zeros in
$(T,2T]$ that are at normalized distance at least $L^2$ from both endpoints.
The number of zeros in $I'\setminus I$ is $O(T^{1/2}L)=o(N)$, and
Lemma~\ref{lem:kernel-limit} below shows that the additional endpoint
deletion costs $O(L^2)=o(N)$.  Pinching $M$ into $M^\circ$ and its complement
gives $\cD(M)\ge\cD(M^\circ)$: pinching is an average of unitary
conjugations, $X\mapsto\tr\Psi(X)$ is convex and unitarily invariant, and
$\Psi\ge0$.

In the notation of~\cite[Equation~(4.4)]{Claude26},
$\widehat A=\widehat G-\widehat E$.  Proposition~4.2 there gives, at
$\lambda=1$,
\[
 \|\widehat E\|_1\ll T^{-1/2},
 \qquad \|\widehat G\|_{\Frob}\ll N^{1/2}.
\]
Consequently
\[
 4\tr\widehat A-\|\widehat A\|_{\Frob}^2
 \ge4\tr\widehat G-\|\widehat G\|_{\Frob}^2-o(N).
\]
The optimized-window form of Theorem~5.8 and the proof of Theorem~D
in~\cite{Claude26} give
\[
 \tr\widehat G=N+o(N),\qquad
 \|\widehat G\|_{\Frob}^2=\frac1{c_1^*}N+o(N),
\]
where $2-1/c_1^*=H_{\MT}$.  Also $N(I')=N+o(N)$.  Substitution in
\eqref{eq:s1-finite}, followed by deletion of the $o(N)$ noncentral points,
yields~\eqref{eq:global-defect}.  Remark~6.1 of~\cite{Claude26} confirms
that the endpoint $\lambda=1$ used here is admissible; its relative source
error is $O(\log\log T/\log T)=o(1)$.
\end{proof}

\section{The Montgomery--Taylor overlap kernel}\label{sec:kernel}

Let $h:=2\pi/L$ and normalize ordinates by
\[
 x_\rho:=\frac{\gamma_\rho-T}{h}
 =\frac{L(\gamma_\rho-T)}{2\pi}.
\]

\begin{lemma}[Uniform finite-Gram comparison]\label{lem:kernel-limit}
After deleting simple zeros in normalized endpoint strips of width $L^2$,
the number deleted is $O(L^2)=o(N)$, and uniformly for all retained simple
zeros,
\begin{equation}\label{eq:kernel-limit}
 \langle v_\rho,v_{\rho'}\rangle
 =k(x_\rho-x_{\rho'})+O(L^{-1}),
\end{equation}
where
\begin{align}
 k(x)&:=\frac{K(x)}{K(0)},\label{eq:k-def}\\
 K(x)&:=\int_{-1/2}^{1/2}
 \cos(\sqrt2t)\cos(2\pi xt)\,dt\notag\\
 &=\frac{\sin(\pi x-1/\sqrt2)}{2\pi x-\sqrt2}
 +\frac{\sin(\pi x+1/\sqrt2)}{2\pi x+\sqrt2},\label{eq:K-closed}\\
 K(0)&=\sqrt2\sin(1/\sqrt2).\notag
\end{align}
The apparent singularities in~\eqref{eq:K-closed} are removable.
\end{lemma}

\begin{proof}
Put
\[
 q_L(t):=\phi_L(Lt)^2,
 \qquad q(t):=\cos(\sqrt2t)\mathbf1_{[-1/2,1/2]}(t).
\]
The two functions differ only on two intervals of total length $2/L$ and
are bounded by one.  Hence
\begin{equation}\label{eq:q-L1}
 \|q_L-q\|_1\le\frac2L,
 \qquad |a_L-K(0)|\le\frac2L.
\end{equation}
For large $L$, $a_L\ge K(0)/2$.  Substitution $u=Lt$ in the full-grid
identity of~\cite[Lemma~2.2]{Claude26}, together with~\eqref{eq:q-L1}, gives
uniformly for every real $x$,
\begin{equation}\label{eq:full-kernel-error}
 \left|
 \frac{\displaystyle\int_{-1/2}^{1/2}
 q_L(t)e^{-2\pi ixt}\,dt}{a_L}-\frac{K(x)}{K(0)}
 \right|\ll L^{-1}.
\end{equation}

It remains to truncate the full grid.  The derivative bounds stated in the
proof of~\cite[Theorem~D]{Claude26} imply, with a constant depending only on
the fixed ramp,
\begin{equation}\label{eq:phi-decay}
 |\widehat\phi_L(r)|\le C_\varrho |r|^{-2}\qquad(r\ne0),
\end{equation}
uniformly in $L$.  For a retained ordinate, the distance to every omitted
grid point is at least
\[
 D_L:=h(L^2-2)\ge\pi L
\]
for all large $L$.  On either side of the grid,
\begin{equation}\label{eq:grid-square-tail}
 \sum_{j\ge0}(D_L+jh)^{-4}
 \le D_L^{-4}+\frac1{3hD_L^3}=O(L^{-2}).
\end{equation}
By Cauchy--Schwarz,~\eqref{eq:phi-decay}, and
\eqref{eq:grid-square-tail}, the omitted product tail is $O(L^{-2})$.
After the Gram normalization by $a_LL^2$ it is $O(L^{-4})$.  Combining this
with~\eqref{eq:full-kernel-error} proves~\eqref{eq:kernel-limit}, uniformly
in the two retained ordinates.

Finally, each deleted normalized strip has physical width
$hL^2=2\pi L$.  The local zero count
$N(t+1)-N(t)\ll\log(t+3)$ gives $O(L\log T)=O(L^2)=o(N)$ deleted zeros.
The elementary integral gives~\eqref{eq:K-closed} and its value at zero.
\end{proof}

Set
\begin{equation}\label{eq:w-def}
 w(x):=k(x)^2.
\end{equation}
The verifier uses the equivalent entire expression
\begin{equation}\label{eq:sinc}
 K(x)=\frac12\left(
 \operatorname{sinc}(\pi x-1/\sqrt2)
 +\operatorname{sinc}(\pi x+1/\sqrt2)\right),
\end{equation}
where $\operatorname{sinc}(z)=\sin z/z$ with its continuous value at zero.

\section{The certified seven-point inequality}\label{sec:certificate}

For six nonnegative consecutive gaps define
\begin{equation}\label{eq:F6}
 \cF_6(g_1,\ldots,g_6)
 :=\frac1{3000}\sum_{i=1}^6g_i
 +\sum_{r=1}^6\frac2{7-r}
 \sum_{i=1}^{7-r}w(g_i+\cdots+g_{i+r-1}).
\end{equation}
The second term contains all 21 pair separations among seven ordered points.

\begin{proposition}[Certified local inequality]\label{prop:F6}
For all $g_1,\ldots,g_6\ge0$,
\begin{equation}\label{eq:F6-bound}
 \cF_6(g_1,\ldots,g_6)\ge C_*
 =\frac{38262312113}{10^{13}}.
\end{equation}
\end{proposition}

\begin{proof}[Computer-assisted verification]
The accompanying verifier uses \texttt{python-flint==0.8.0} and Arb interval
arithmetic~\cite{Johansson17}.  Its grid mesh is $1/4000$.  On every closed
cell it evaluates~\eqref{eq:sinc}, forms a rigorous lower bound for $w$, and
widens conversion to IEEE-754 binary64 downward.  A separate interval table
lower-bounds $w''$.  All later range queries use minima of those cell bounds.

The exact pressure cutoff is
\[
 q=\left\lceil C_*\cdot4000\cdot3000\right\rceil=45915.
\]
Thus $\sum_i g_i/3000$ proves the target whenever the sum of lower cell
indices is at least $q$.  Before forming six-dimensional boxes, the verifier
also uses the one-variable contribution
\[
 U(g):=\frac g{3000}+\frac13w(g).
\]
The surviving cells form the three exact components
\[
 [3807,4780],\qquad[7218,9373],\qquad[10560,44945].
\]
Pressure filtering leaves 435 initial products of components.

For each box, every consecutive partial sum in~\eqref{eq:F6} receives a
range-minimum lower bound.  If this direct interval bound is insufficient,
the code lower-bounds the Hessian throughout the box using the $w''$ table.
A binary64 LDL factorization is only a cheap filter; positive definiteness is
rechecked with Arb.  On a certified convex box, evaluation at its rational
midpoint gives the rigorous tangent bound
\[
 \cF_6(x)\ge\cF_6(c)
 -\sum_{i=1}^6|\partial_i\cF_6(c)|r_i.
\]

Two neighborhoods of the numerical minimizers require a sharper local
argument.  The first has exact rational center represented by
\begin{equation}\label{eq:basin-center}
 \begin{aligned}
 c=(&1.04608035577,\,1.98913202062,\,1.98641493611,\\
    &1.04160329372,\,1.97702352234,\,1.04500209462).
 \end{aligned}
\end{equation}
and the second has reversed coordinates.  Each coordinate radius is exactly
$1/64$.  Interval lower bounds for $w''$, followed by an Arb LDL check, prove
throughout both basins that
\begin{equation}\label{eq:strong-convexity}
 \nabla^2\cF_6\succeq\frac3{16}I.
\end{equation}
At 256-bit basin precision, strong convexity gives
\begin{align}
 \inf_{\text{basin}}\cF_6
 &\ge \cF_6(c)-\frac{\|\nabla\cF_6(c)\|^2}{2(3/16)}\notag\\
 &\ge0.00382623121130447424827371713006270699383823035,
 \label{eq:basin-lower}
\end{align}
and the same lower endpoint for the reversed basin.  This clears the exact
target by more than $4.47\times10^{-15}$.  A search box is basin-pruned only
when all of its closed grid cells lie inside a certified basin.

All other boxes are bisected along a widest coordinate.  Reaching an
unresolved terminal cell raises an exception.  The canonical run has the
following deterministic counts.
\begin{center}
\begin{tabular}{lr}
\toprule
Field & Canonical value\\
\midrule
Initial boxes & 435\\
Visited nodes & 822,433\\
Pruned leaves & 411,434\\
Splits & 410,999\\
Maximum depth & 39\\
Interval prunes & 285,258\\
Pressure prunes & 2,872\\
Tangent prunes & 122,329\\
Strong-convexity prunes & 975\\
Unresolved terminal cells & 0\\
\bottomrule
\end{tabular}
\end{center}
The code checks both binary-forest identities
\[
 822433=411434+410999,
 \qquad411434=410999+435.
\]
The canonical kernel-table SHA-256 is
\begin{sloppypar}
\small\ttfamily
5a1f95f754a83ba05f37692d4bda69bf3fa5d3752af902826bfc4cffd31428b9,
\end{sloppypar}
and the second-derivative-table SHA-256 is
\begin{sloppypar}
\small\ttfamily
318ee79f619cc75c6f7a30424a181764865910693f33a3d95469ebcae0ce42ba.
\end{sloppypar}

A full rerun with table precision increased from 128 to 160 bits and basin
precision from 256 to 320 bits produces the same initial boxes, node counts,
pruning counts, and depth.  Its independently regenerated second-derivative
table has SHA-256
\begin{sloppypar}
\small\ttfamily
fd634c87441380105aad52f069f735108d0dabcd7d5d0a993b6abc52ee12b055.
\end{sloppypar}
This is a cross-precision stress test, not an independent implementation: it
uses the same source and Arb/FLINT library.  The exact expected reports are
included in the artifact and are compared field by field on every release
run.  At commit
\begin{center}
\small\ttfamily b8dcb42c88f920fca2c609b1c88a5cc133d88b2a,
\end{center}
clean local replays with the hash-locked CPython~3.12 wheel ended respectively
with
\begin{center}
\small\ttfamily
release\_certificate=match,\qquad
cross\_precision\_certificate=match.
\end{center}

The additional script
\texttt{verifier/scripts/verify\_basin\_independent.py} does not import the
release verifier.  It reimplements the sinc quotient derivatives,
subdivides each scalar basin interval into 2,048 exact rational pieces, and
performs the shifted $6\times6$ LDL test over exact rational arithmetic.  It
proves~\eqref{eq:strong-convexity} for both basins and obtains the lower
endpoint
\[
 0.003826231211304474248273717130062706993838230354884912386.
\]
This is a separate control path but retains Arb/FLINT as the transcendental
interval trust base.  These facts complete the finite verification of
Proposition~\ref{prop:F6}.
\end{proof}

\begin{remark}
The center~\eqref{eq:basin-center} was located by ordinary floating-point
optimization, but the optimizer output is not trusted.  Only the displayed
terminating decimals, interpreted as exact rationals, and the interval proof
of~\eqref{eq:strong-convexity}--\eqref{eq:basin-lower} enter the certificate.
\end{remark}

\section{From seven-point windows to block defect}\label{sec:block}

Let $y_1<\cdots<y_m$ and set
\begin{equation}\label{eq:Em}
 E_m:=2\sum_{1\le i<j\le m}w(y_j-y_i).
\end{equation}

\begin{lemma}[Block energy]\label{lem:block-energy}
For every $m\ge7$,
\begin{equation}\label{eq:block-energy}
 E_m+\frac1{500}(y_m-y_1)\ge C_*(m-6).
\end{equation}
\end{lemma}

\begin{proof}
Write $g_i=y_{i+1}-y_i$ and sum~\eqref{eq:F6-bound} over the $m-6$
consecutive seven-point windows.  A pair spanning $r\le6$ gaps occurs in at
most $7-r$ windows and has coefficient $2/(7-r)$ in each, so its total
coefficient is at most $2$.  Pairs spanning more than six gaps do not appear
and contribute nonnegatively to $E_m$.  Each single gap occurs in at most six
windows, so its total linear coefficient is at most $6/3000=1/500$.
\end{proof}

\begin{lemma}[Block defect]\label{lem:block-defect}
If $G\succeq0$ is Hermitian, then
\begin{equation}\label{eq:block-defect}
 \tr\Psi(G)\ge
 \min\left\{1,2\sum_{i<j}|G_{ij}|^2\right\}.
\end{equation}
\end{lemma}

\begin{proof}
If every eigenvalue of $G$ is at most $2$, then
$\Psi(G)=(G-I)^2$ and
\[
 \tr\Psi(G)=\tr(G-I)^2
 \ge2\sum_{i<j}|G_{ij}|^2.
\]
If some eigenvalue $\lambda>2$, then $\Psi(\lambda)=2\lambda-3>1$, while
all other contributions are nonnegative.
\end{proof}

Choose
\begin{equation}\label{eq:m-A0}
 m:=267,
 \qquad A_0:=C_*(m-6)=0.9986463461493<1.
\end{equation}
For comparison,
\begin{equation}\label{eq:next-m}
 C_*(268-6)=1.0024725773606>1.
\end{equation}
Thus $m=267$ is the last block size before the unit cap in
Lemma~\ref{lem:block-defect} activates.  More explicitly, for general $m$
the block constant is $A_m=\min\{1,C_*(m-6)\}$ and the same argument gives
\[
 R_m=\frac{H_{\MT}-(m-1)/(500m)}{1-A_m/m}.
\]
On $7\le m\le267$ this is increasing because
$6C_*(H_{\MT}-1/500)>(1-C_*)/500$, while for $m\ge268$ it is decreasing
with derivative $-H_{\MT}/(m-1)^2$ after writing
$R_m=((H_{\MT}-1/500)m+1/500)/(m-1)$.  Hence $m=267$ maximizes the
resulting bound among integer block sizes $m\ge7$.

Let $B$ be a block of 267 consecutive retained central simple zeros, and let
$G_B$ be its Gram matrix.  Lemma~\ref{lem:kernel-limit} and the bounds
$|(G_B)_{ij}|,|k(x_i-x_j)|\le1$ give, uniformly in $B$,
\begin{equation}\label{eq:block-energy-comparison}
 \left|2\sum_{i<j}|(G_B)_{ij}|^2-E_m\right|
 \le\frac{C_m}{L},
\end{equation}
where $C_m$ is absolute because $m=267$ is fixed.  Put
$Q_B=2\sum_{i<j}|(G_B)_{ij}|^2$.  If $Q_B\le1$, then
Lemmas~\ref{lem:block-energy}--\ref{lem:block-defect} and
\eqref{eq:block-energy-comparison} give the next inequality.  If $Q_B>1$,
then $\cD(G_B)\ge1>A_0$, so it holds again.  Thus, uniformly in $B$,
\begin{equation}\label{eq:267-block}
 \cD(G_B)+\frac1{500}\Span(B)\ge A_0-O(L^{-1}).
\end{equation}

\section{Shifted block pinching and final arithmetic}\label{sec:pinching}

Let
\[
 x_1<\cdots<x_{S^\circ}
\]
be the normalized ordinates of the retained central simple zeros.  Endpoint
deletion gives
\begin{equation}\label{eq:S-central}
 S^\circ=S-O(L^2)=S-o(N).
\end{equation}
For each of the $m=267$ offsets, partition the ordered points into full
consecutive blocks of size $m$, leaving only bounded endpoint remainders.
Pinching $M^\circ$ to the principal blocks and the remainder gives
\begin{equation}\label{eq:block-pinching}
 \cD(M^\circ)\ge\sum_B\cD(G_B).
\end{equation}

For each offset, sum~\eqref{eq:267-block} over its full blocks and use
\eqref{eq:block-pinching}; then average over all $m$ offsets.  The average
number of full blocks is $S^\circ/m+O(1)$.  Each adjacent gap belongs to a
block span for at most $m-1$ offsets.  Moreover
\begin{equation}\label{eq:total-span}
 x_{S^\circ}-x_1
 \le\frac{LT}{2\pi}=N+o(N)
\end{equation}
by the Riemann--von Mangoldt formula.  Across all offsets there are at most
$S^\circ\le N$ full blocks, each occurring once, so the total error from
\eqref{eq:267-block} is $O(N/L)=o(N)$.  Consequently
\begin{equation}\label{eq:defect-global}
 \cD(M^\circ)
 \ge\frac{A_0}{m}S-\frac{m-1}{500m}N-o(N)
 =\frac{261C_*}{267}S-\frac{266}{133500}N-o(N).
\end{equation}

Substitute~\eqref{eq:defect-global} into~\eqref{eq:global-defect}:
\[
 S\ge H_{\MT}N+\frac{261C_*}{267}S
 -\frac{266}{133500}N-o(N).
\]
Hence
\begin{equation}\label{eq:solve}
 \left(1-\frac{261C_*}{267}\right)S
 \ge\left(H_{\MT}-\frac{266}{133500}\right)N-o(N).
\end{equation}
Divide by $N$, let $T\to\infty$, and use~\eqref{eq:Cstar}.  This gives
\eqref{eq:main-formula}.  Multiplying numerator and denominator by
$267\cdot10^{13}$ gives the exact reduced expression in
Theorem~\ref{thm:main}; the common integer gcd is one.  This proves
Theorem~\ref{thm:main}.

\section{Position relative to prior work}\label{sec:prior}

Montgomery's pair-correlation argument proved, under RH, that at least two
thirds of the zeros are simple~\cite{Montgomery73}; the optimized
Montgomery--Taylor window gives the constant~\eqref{eq:HMT}
\cite{Montgomery75}.  Chirre, Gon\c{c}alves, and de Laat later obtained
stronger conditional simple-zero estimates by semidefinite programming and
positivity outside the bandwidth-one interval~\cite{CGdL20}.  Those results
operate in a different conditional regime and are not numerically comparable
to Theorem~\ref{thm:main} as an unconditional critical-line statement.

Recent work of Goldston and collaborators clarified how pair correlation can
control zeros that are simultaneously simple and on the critical line, under
the pair-correlation conjecture or narrow-box hypotheses
\cite{GLSS25,GS25,GS26}.  Claude's paper~\cite{Claude26} then supplied the
unconditional finite-dimensional inertia framework and the baseline
$H_{\MT}$ used here.  The present refinement does not extend Fourier support
or add a new prime-side estimate.  It extracts extra information from the
short-range Gram matrix by preserving the spectral defect $\cD$ and
certifying a seven-point local inequality.

The software lineage is also relevant.  A public AI-generated artifact first
implemented the stability approach at local target $0.0038$
\cite{Ainta26}; a clean-room release certified $0.00382$
\cite{Rademacher26}.  The current artifact certifies~\eqref{eq:Cstar} and adds
the explicit strong-convexity basins and precision replay.  These repositories
are research artifacts, not substitutes for peer-reviewed priority analysis.
A search of the cited 2025--2026 primary literature found no identical
stability-defect/seven-point construction, but that search cannot establish
priority.

\section{Reproducibility, trust base, and limitations}
\label{sec:reproducibility}

Proposition~\ref{prop:F6} is the only computer-assisted proposition in the
deduction.  The accompanying source archive contains the verifier, exact JSON
reports, unit tests, a hash-locked dependency file, and detailed audit and
reproduction notes.  The canonical commands are
\begin{verbatim}
python3.12 -m venv .venv
source .venv/bin/activate
python -m pip install --require-hashes -r requirements.lock
PYTHONPATH=src python -m unittest discover -s tests -v
PYTHONPATH=src python scripts/verify_release.py
PYTHONPATH=src python scripts/verify_cross_precision.py
python scripts/verify_basin_independent.py
\end{verbatim}
The two verifier scripts agree with their expected reports when they end with
the respective status strings
\begin{center}
\small\ttfamily
release\_certificate=match\\[2pt]
cross\_precision\_certificate=match.
\end{center}

The partial Lean companion uses Lean \texttt{v4.33.0-rc2} and is pinned to
the following upstream and mathlib commits:
\begin{center}
\small\ttfamily
upstream: 3635e74826a4c1fcece7d1cd2b6fa75e43a00510\\[2pt]
mathlib: 51e6992efd06126df61a496bebf8f49482a4e129.
\end{center}
Its audit commands are
\begin{verbatim}
lake build ZetaSeven
lake env lean ZetaSeven/Audit.lean
\end{verbatim}
The proved declarations cover the rank--trace surplus, the simple-zero Gram
split and its inertia bound, the finite tail seam, the bridge from
off-diagonal Gram energy to spectral defect, the exact consecutive-window
count, arbitrary principal-block spectral pinching, abstract 267-offset
aggregation, the conditional 267-point block bridge, the abstract
defect-preserving Theorem~D source inequality (including an explicit
$o(N)$ error package), the concrete shifted partitions and endpoint fibers,
their insertion into the exact finite simple-zero count, and the exact
rational and epsilon-form assembly.
The proposition \texttt{SevenPointClaim} is intentionally defined but not
asserted as a theorem.  Consequently, a successful Lean build is evidence
for these layers only; it is not an end-to-end formal proof of
Theorem~\ref{thm:main}.

The finite trust base consists of CPython 3.12, IEEE-754 binary64 semantics,
\texttt{python-flint} version 0.8.0, Arb/FLINT, the operating system and hardware,
and the included source.  The verifier rebuilds all transcendental tables; it
does not trust cached tables, the numerical optimizer, or committed run logs.
The separate basin checker changes the derivative implementation,
subdivision, and exact linear-algebra path, but uses the same Arb/FLINT
library.  No fully independent interval library or proof-assistant checker is
yet available for the complete six-variable subdivision.  Lean reduces the
trust base for the surrounding algebra but does not reduce the Arb
certificate's trust base until a proof-carrying certificate checker is
completed.

The larger mathematical trust boundary is equally important.  This work
does not independently reprove the prime-side Theorem~5.8 of~\cite{Claude26}
or its explicit-formula inputs.  Corollary~\ref{cor:global-defect} and
Lemma~\ref{lem:kernel-limit} spell out the new interface, endpoint deletion,
finite-grid tail, and uniform Gram comparison rather than leaving them as
formal hypotheses.  Reviewers should nevertheless audit their normalization
against the upstream concrete parameter family and independently assess the
still-new cited theorem.

For peer review, the highest-value checks are:
\begin{enumerate}[leftmargin=*,itemsep=2pt]
 \item verify Lemma~\ref{lem:stable-rank-trace} and the deduction of
 Corollary~\ref{cor:global-defect} against the exact notation of
 \cite{Claude26};
 \item independently inspect the cell-covering, outward-rounding, Hessian,
 and basin-containment logic in the verifier;
 \item reproduce both certificates on a clean machine and, preferably, a
 second architecture;
 \item implement an independent interval or proof-assistant checker for
 Proposition~\ref{prop:F6};
 \item check the seven-window multiplicities and shifted gap charges in
 Sections~\ref{sec:block}--\ref{sec:pinching};
 \item repeat the prior-art search before any public priority claim.
\end{enumerate}

Finally, neither a positive proportion result nor the local statistics used
here imply that all zeros lie on the critical line.  The method remains far
from a proof of the Riemann hypothesis, and the bandwidth-one framework has
known structural ceilings discussed in~\cite{Claude26}.

\section*{Acknowledgment and authorship placeholder}

This anonymous line is intended for review circulation only.  A submission
must name the people who take responsibility for the mathematics, disclose
the AI-assisted origin in accordance with the venue's policy, and obtain
their approval of the final text.  The upstream software credits and license
are recorded in the accompanying \texttt{PROVENANCE.md}.

\begin{thebibliography}{99}

\bibitem{Claude26}
Claude,
\emph{More than two thirds of the zeros of the Riemann zeta function lie on
the critical line}, preprint, August 10, 2026,
\url{https://www-cdn.anthropic.com/564f962e60643842f5fcb4a17c9dbc8f608f1c37.pdf}.

\bibitem{ClaudeLean26}
Anthropic,
\emph{Zeta23: Lean 4 formalization of ``More than two thirds of the zeros of
the Riemann zeta function lie on the critical line''}, 2026,
\url{https://github.com/anthropics/zeta-23-lean}.

\bibitem{Johansson17}
F.~Johansson,
\emph{Arb: efficient arbitrary-precision midpoint-radius interval
arithmetic}, IEEE Trans. Comput. \textbf{66} (2017), no.~8, 1281--1292.

\bibitem{Montgomery73}
H.~L.~Montgomery,
\emph{The pair correlation of zeros of the zeta function}, in: Analytic
Number Theory (St. Louis, 1972), Proc. Sympos. Pure Math. \textbf{24},
Amer. Math. Soc., 1973, 181--193.

\bibitem{Montgomery75}
H.~L.~Montgomery,
\emph{Distribution of the zeros of the Riemann zeta function}, in:
Proceedings of the International Congress of Mathematicians (Vancouver,
1974), Vol.~1, Canad. Math. Congress, 1975, 379--381.

\bibitem{CGdL20}
A.~Chirre, F.~Gon\c{c}alves, and D.~de Laat,
\emph{Pair correlation estimates for the zeros of the zeta function via
semidefinite programming}, Adv. Math. \textbf{361} (2020), 106926;
\url{https://arxiv.org/abs/1810.08843}.

\bibitem{GLSS25}
D.~A.~Goldston, J.~Lee, J.~Schettler, and A.~I.~Suriajaya,
\emph{Pair correlation conjecture for the zeros of the Riemann zeta-function
I: simple and critical zeros}, arXiv:2503.15449 (2025),
\url{https://arxiv.org/abs/2503.15449}.

\bibitem{GS25}
D.~A.~Goldston and A.~I.~Suriajaya,
\emph{Zeta zeros on the critical line}, arXiv:2511.20059 (2025),
\url{https://arxiv.org/abs/2511.20059}.

\bibitem{GS26}
D.~A.~Goldston and A.~I.~Suriajaya,
\emph{Zeta zeros in a narrow vertical box}, arXiv:2603.28104 (2026),
\url{https://arxiv.org/abs/2603.28104}.

\bibitem{Ainta26}
ainta,
\emph{zeta-simple-zeros: reproducible verification for a simple zeta-zero
bound}, software artifact, commit
\texttt{040c5e899e658aed7b56a2a87f501798fe10761d}, 2026,
\url{https://github.com/ainta/zeta-simple-zeros}.

\bibitem{Rademacher26}
L.~Rademacher,
\emph{AI refines AI: a reproducible candidate improvement to the simple
zeta-zero bound}, software artifact, commit
\texttt{bd4a7d36988b23220034527881f4457f2f689e86}, 2026,
\url{https://github.com/learademacher/ai-refines-ai-zeta-bound}.

\end{thebibliography}

\end{document}
